{\color{ink}\bfseries 例 7-7}\par\vspace{1bp}本题按\(1 MB=\allowbreak{}10^{6} B\)计算，32位缓冲可存4 B。\par （1）A的取数间隔上限为\(4/\allowbreak{}(2\times 10^{6})=\allowbreak{}2 \mu s\)；每次至少\(10\times 4=\allowbreak{}40\)周期，即\(40/\allowbreak{}(500\times 10^{6})=\allowbreak{}80 ns\)。每秒需服务\(5\times 10^{5}\)次，CPU占用至少\((5\times 10^{5})\times 80\times 10^{-9}=\allowbreak{}4\%\)。\par （2）B每\(4/\allowbreak{}(40\times 10^{6})=\allowbreak{}0.1 \mu s\)到达一缓冲数据，每次中断却至少占\(400/\allowbreak{}(500\times 10^{6})=\allowbreak{}0.8 \mu s\)。所需CPU比例为\(0.8/\allowbreak{}0.1=\allowbreak{}8\)，即800\%，不能按此粒度采用中断I/\allowbreak{}O。\par （3）DMA每秒最多完成\(40\times 10^{6}/\allowbreak{}1000=\allowbreak{}4\times 10^{4}\)块，每块预后处理占\(500/\allowbreak{}(500\times 10^{6})=\allowbreak{}1 \mu s\)，CPU占用最多\((4\times 10^{4})\times 10^{-6}=\allowbreak{}4\%\)。软件开销按1000 B块计数，接口的4 B决定块内传送粒度。
